\documentclass[12 pt, letterpaper]{hmcpset}
\usepackage{hw-preamble}

\name{Jayden Thadani}
\class{Mathematical Analysis I --- MATH 131}
\assignment{Homework 6}
\duedate{12th October 2024}

\begin{document}

Collaborators: Mithra Karamchedu

\begin{problem}[Rudin chapter 3, exercise 6]
	Investigate the behaviour (convergence or divergence) of $\sum a_{n}$ if
	\begin{enumerate}
		\item $a_{n} = \sqrt{n + 1} - \sqrt{n}$
		\item $a_{n} = \frac{\sqrt{n + 1} - \sqrt{n}}{n}$ 
		\item $a_{n} = (\sqrt[n]{n} - 1)^{n}$
		\item $a_{n} = \frac{1}{1 + z^{n}}$ for complex $z$.
	\end{enumerate}
\end{problem}

\begin{solution}
	\begin{enumerate}
		\item
			\[
				\boxed{
					\sum_{n = 1}^{\infty} \sqrt{n + 1} - \sqrt{n} \text{ diverges to } \infty.
				}
			\]

			Consider the partial sums
			\[
				\begin{split}
					s_{N} & = \sum_{n = 1}^{N} a_{n} \\
					      & = (\sqrt{N + 1} - \sqrt{N}) + (\sqrt{N} - \sqrt{N - 1}) + \dots + (\sqrt{2} - \sqrt{1}) \\
					      & = \sqrt{N + 1} + (\sqrt{N} - \sqrt{N}) + \dots + (\sqrt{2} - \sqrt{2}) - \sqrt{1} \\
					      & = \sqrt{N + 1} - 1.
				\end{split}
			\]
			Note that for any $M > 0$, $n > (M + 1)^{2}$ gives
			\[
				s_{n} = \sqrt{(M + 1)^{2} + 1} - 1 > M + 1 - 1 = M
			\]
			so $s_{n} \to \infty$ and thus, $\sum a_{n} = \infty$.
		\item
			\[
				\boxed{
					\sum_{n = 1}^{\infty} \frac{\sqrt{n + 1} - \sqrt{n}}{n} = 0.
				}
			\]

			Again, we consider the partial sums
			\[
				\begin{split}
					s_{N} & = \sum_{n = 1}^{N} a_{n} \\
					      & = \sum_{n = 1}^{N} \frac{\sqrt{n + 1} - \sqrt{n}}{n} \\
					      & < \sum_{n = 1}^{N} \frac{\sqrt{n + 2 \sqrt{n} + 1} - \sqrt{n}}{n} \\
					      & = \sum_{n = 1}^{N} \frac{\sqrt{n} - \sqrt{n} + 1}{n} \\
					      & = \frac{1}{N}.
				\end{split}
			\]

			By the Archimedean property, $s_{n} \to 0$ and thus, $\sum a_{n} = 0$.
		\item
			\[
				\boxed{
					\sum_{n = 1}^{\infty} (\sqrt[n]{n} + 1)^{n} \text{ converges.}
				}
			\]

			By the root test, $\sum a_{n}$ converges since
			\[
				\begin{split}
					\limsup_{n \to \infty} \lvert a_{n} \rvert^{1/n} & = \lim_{n \to \infty} (n^{1/n} - 1) \\
											 & = (\lim_{n \to \infty} n^{1/n}) - 1 \\
											 & = 0 \\
											 & < 1.
				\end{split}
			\]
		\item
			\[
				\boxed{
					\sum_{n = 1}^{\infty} \frac{1}{1 + z^{n}} \text{ converges for } \lvert z \rvert < 1 \text{ and diverges otherwise.}
				}
			\]

			We compare $\lvert 1 + z^{n} \rvert$ to $n^{p}$ for some fixed $p > 1$. If $\lvert z \rvert > 1$, then for large enough $n$ we have $\lvert 1 + z^{n} \rvert > n^{p}$ (this follows from Theorem 3.20). Thus,
			\[
				\sum_{n = 1}^{\infty} \frac{1}{\lvert 1 + z^{n} \rvert} = \sum_{n = 1}^{\infty} \frac{1}{n^{p}}
			\]
			and converges. It follows from the absolute convergence that the series itself converges.

			However, if $\lvert z \rvert \le 1$ then for $n > 2$ we have $\lvert 1 + z^{n} \rvert < n$. Thus,
			\[
				\sum_{n = 1}^{\infty} \frac{1}{n} < \sum_{n = 1}^{\infty} \frac{1}{\lvert 1 + z^{n} \rvert}
			\]
			and since the harmonic series diverges, so does the series here.
	\end{enumerate}
\end{solution}

\pagebreak

\begin{problem}[Rudin chapter 3, exercise 7]
	Prove that the convergence of $\sum a_{n}$ implies the convergence of
	\[
		\sum \frac{\sqrt{a_{n}}}{n}
	\]
	if $a_{n} \ge 0$.
\end{problem}

\begin{solution}
	By comparison with harmonic series, we know that since $a_{n}$ converges, it cannot exceed the harmonic series at infinitely many $n$. That is, for sufficiently large $n$, we should have $a_{n} < 1/n$. But then $\sqrt{a_{n}} < 1/\sqrt{n}$ and thus, $\sqrt{a_{n}}/n < 1/n^{3/2}$.

	We know $\sum 1/n^{3/2}$ converges. Thus, by the comparison test we also have that $\sum \sqrt{a_{n}}/n$ converges.
\end{solution}

\pagebreak

\begin{problem}[Rudin chapter 3, exercise 8]
	If $\sum a_{n}$ converges and if $\{ b_{n} \}_{n}$ is monotonic and bounded, prove that $\sum a_{n} b_{n}$ converges.
\end{problem}

\begin{solution}
	Since $b_{n}$ is monotonic and bounded it must converge to some $b$. Given $\epsilon > 0$, choose $ \varepsilon_{0} < \min \{ \varepsilon/2b, \sqrt{\varepsilon}/2 \}$ Suppose for all $m > N_{b}$, we have  $\lvert b_{m} - b \rvert < \varepsilon_{0}$. Also suppose that for all $n > m > N_{a}$ we have $\left \lvert \sum_{k = m + 1}^{n} a_{k} \right \rvert < \varepsilon_{0}$.

	Then for all $n, m > \max \{ N_{a}, N_{b} \}$ we have
	\[
		\begin{split}
			\left \lvert \sum_{k = m + 1}^{n} a_{k} b_{k} \right \rvert & = \left \lvert \sum_{k = m + 1}^{n} a_{k} b + \sum_{k = m + 1}^{n} a_{k} (b_{k} - b) \right \rvert \\
										    & \le \lvert b \rvert \left \lvert \sum_{k = m + 1}^{n} a_{k} \right \rvert + \lvert b_{k} - b \rvert \left \lvert \sum_{k = m + 1}^{n} a_{k} \right \rvert \\
										    & < \lvert b \rvert \varepsilon_{0} + \varepsilon_{0}^{2} \\
										    & < \varepsilon/2 + \varepsilon/2 \\
										    & = \varepsilon.
		\end{split}
	\]

	That is, the partial sums of $\sum a_{n} b_{n}$ form a Cauchy sequence. Thus, $\sum a_{n} b_{n}$ must converge.
\end{solution}

\pagebreak

\begin{problem}[Rudin chapter 3, exercise 9]
	Find the radius of convergence of each of the following power series:
	\begin{enumerate}
		\item $\sum n^{3} z^{n}$ 
		\item $\sum \frac{2^{n}}{n!} z^{n}$ 
		\item $\sum \frac{2^{n}}{n^{2}} z^{n}$ 
		\item $\sum \frac{n^{3}}{3^{n}} z^{n}$
	\end{enumerate}
\end{problem}

\begin{solution}
	\begin{enumerate}
		\item
			\[
				\boxed{
					R = 1.
				}
			\]

			Let $\alpha = \limsup_{n \to \infty} \lvert n^{3} \rvert^{1/n}$ which can be written as $\lim_{n \to \infty} (n^{1/n})^{3}$. Then $R = 1/\alpha$. We have shown in class that $\lim_{n \to \infty} n^{1/n} = 1$. Thus,
			\[
				\lim_{n \to \infty} (n^{1/n})^{3} = (\lim_{n \to \infty} n^{1/n})^{3} = 1. 
			\]
			But then $\alpha = \lim_{n \to \infty} (n^{1/n})^{3} = 1$ and thus, $R = 1$.
		\item
			\[
				\boxed{
					R = \infty.
				}
			\]

			Let $\alpha = \limsup \left \lvert \frac{2^{n}}{n!} \right \rvert^{1/n}$. Then $R = 1/\alpha$. We claim that $(n!)^{1/n} \ge \sqrt{n}$ for sufficiently large $n$, and thus,
			\[
				\left \lvert \frac{2^{n}}{n!} \right \rvert^{1/n} \le \frac{2}{\sqrt{n}}
			\]
			We can show this by comparing $(n!)^{2}$ and $n^{n}$. If we write
			\[
				(n!)^{2} = n (2(n - 1)) (3(n - 2)) \dots ((n - 2) 3) ((n - 1) 2) n
			\]
			we can see that for sufficiently large $n$, each factor is clearly greater than $n$. Thus $(n!)^{2} > n^{n}$.

			Now we have
			\[
				\alpha \le \lim_{n \to \infty} \frac{2}{\sqrt{n}} = 0
			\]
			So $\alpha = 0$, and the series must converge on all of $\mathbb{C}$. That is, $R = \infty$.
		\item
			\[
				\boxed{
					R = \frac{1}{2}.
				}
			\]

			Let $\alpha = \limsup \left \lvert \frac{2^{n}}{n^{2}} \right \rvert^{1/n}$. Notice that
			\[
				\begin{split}
					\lim_{n \to \infty} \left \lvert \frac{2^{n}}{n^{2}} \right \rvert^{1/n} & = \lim_{n \to \infty} \frac{2}{n^{2/n}} \\
														 & = \frac{2}{(\lim_{n \to \infty} n^{1/n})^{2}} \\
														 & = 2.
				\end{split}
			\]

			Since the limit exists, the $\limsup$, $\alpha$ must be the limit. That is, $\alpha = 2$ and $R = 1/2$.
		\item
			\[
				\boxed{
					R = 3.
				}
			\]

			Let $\alpha = \limsup_{n \to \infty} \lvert n^{3}/3^{n} \rvert^{1/n}$. Thus, $R = 1/\alpha$. But we've already shown that
			\[
				\limsup_{n \to \infty} \lvert n^{3} \rvert^{1/n} = \lim_{n \to \infty} \lvert n^{3} \rvert^{1/n} = 1.
			\]
			Thus,
			\[
				\limsup_{n \to \infty} \left \lvert \frac{n^{3}}{3^{n}} \right \rvert^{1/n} = \lim_{n \to \infty} \frac{\lvert n^{3} \rvert^{1/n}}{3} = \frac{1}{3}.
			\]
			
			Then $R = 3$.
	\end{enumerate}
\end{solution}

\pagebreak

\begin{problem}[Rudin chapter 3, exercise 11]
	Suppose $a_{n} > 0$, $s_{n} = a_{1} + \dots + a_{n}$, and $\sum a_{n}$ diverges.
	\begin{enumerate}
		\item Prove that $\sum \frac{a_{n}}{1 + a_{n}}$ diverges.
		\item Prove that
			\[
				\frac{a_{N + 1}}{s_{N + 1}} + \dots + \frac{a_{N + k}}{s_{N + k}} \ge 1 - \frac{s_{N}}{s_{N + k}}
			\]
			and deduce that $\sum \frac{a_{n}}{s_{n}}$ diverges.
		\item Prove that
			\[
				\frac{a_{n}}{s_{n}^{2}} \le \frac{1}{s_{n - 1}} - \frac{1}{s_{n}}
			\]
			and deduce that $\sum \frac{a_{n}}{s_{n}^{2}}$ converges.
		\item What can be said about
			\[
				\sum \frac{a_{n}}{1 + n a_{n}} \quad \text{and} \quad \sum \frac{a_{n}}{1 + n^{2} a_{n}}?
			\]
	\end{enumerate}
\end{problem}

\begin{solution}
	\begin{enumerate}
		\item If $\sum a_{n}$ then $a_{n}$ must be unbounded. Thus, there exists some $N$ such that for $n > N$, $a_{n} > 1$ is always true. But then
			\[
				\frac{a_{n}}{a_{n} + 1} > \frac{1}{2}
			\]
			is true for all $n > N$. Thus, the summand doesn't go to zero and the sum must diverge.
		\item Note that since $a_{n} > 0$, $\{ s_{n} \}_{n}$ is increasing. Thus
			\[
				\begin{split}
					\frac{a_{N + 1}}{s_{N + 1}} + \dots + \frac{a_{N + k}}{s_{N + k}} & \ge \frac{a_{N + 1} + \dots + a_{N + k}}{s_{N + k}} \\
					& = \frac{s_{N + k} - s_{N}}{s_{N}} \\
					& = 1 - \frac{s_{N}}{s_{N + k}}.
				\end{split}
			\]

			Since $\sum a_{n}$ diverges, its partial sums cannot form a Cauchy sequence. Thus, there is some $k$ with $\lim_{N \to \infty} s_{N}/s_{N + k} \neq 0$. But then
			\[
				\frac{a_{N + 1}}{s_{N + 1}} + \dots + \frac{a_{N + k}}{s_{N + k}} > 1 - s > 0.
			\]
			This exactly implies that the partial sums of $\sum a_{n} / s_{n}$ do not form a Cauchy sequence --- there is no $N$ such that $n > N$ gives us that the $n$th and $n + k$th partial sums are less than $1 - s$ apart. Thus, $\sum a_{n} / s_{n}$ diverges.
		\item Again, recall that  $\{ s_{n} \}_{n}$ is increasing. Thus,
			\[
				\begin{split}
					\frac{1}{s_{n - 1}} - \frac{1}{s_{n}} & = \frac{s_{n} - s_{n - 1}}{s_{n} s_{n - 1}} \\
									      & \ge \frac{a_{n}}{s_{n^{2}}}.
				\end{split}
			\]

			Now we can calculate the partial sums of $\sum a_{n} / s_{n}^{2}$ ---
			\[
					\sum_{n = 1}^{N} \frac{a_{n}}{s_{n}^{2}} = \frac{a_{1}}{s_{1}^{2}} + \frac{1}{s_{1}} - \frac{1}{s_{N}}.
			\]
			Since $s_{N} \to \infty$, we have $1/s_{N} \to 0$, and thus,
			\[
				\sum_{n = 1}^{\infty} \frac{a_{n}}{s_{n}^{2}} = \frac{a_{1}}{s_{1}^{2}} + \frac{1}{s_{1}}.
			\]
		\item Note that
			\[
				\frac{a_{n}}{1 + n^{2} a_{n}} = \frac{1}{1/a_{n} + n^{2}} < \frac{1}{n^{2}}.
			\]
			Thus, by the comparison test, $\frac{a_{n}}{1 + n^{2} a_{n}}$ converges for positive $a_{n}$.
	\end{enumerate}
\end{solution}

\end{document}
